% English translation of checked_pages/esq_119_046.tex; source: cote 119, site 46, PDF page 47, typescript page 44.
% AI-translated by GPT-6 Astra (OpenAI, ChatGPT Pro mode).
\scanpage{119/46}
\mspage{44}

different direction from the one that was and is close to my heart. It certainly
introduces derived categories in various noncommutative
contexts, but without entering into the question of the essential internal
structures of such a category, also left open
in the commutative case by Verdier and, after him, by Illusie. Likewise,
the question of identifying the natural ``coefficients''
for a noncommutative cohomological formalism, beyond
the stacks (which should be called 1-stacks) studied in Giraud's
book, remained open --- or rather, the rich and
precise intuitions that answer it, drawn from many examples
coming in particular from algebraic geometry, still await
a precise and flexible language to give them form.

I return to certain aspects of these foundational
questions in 1975, on the occasion (if I remember correctly) of a
correspondence with Larry Breen (two letters from this correspondence
will be reproduced as an appendix to Chap. I of volume 1,
``History of Models'',\tnote{TN-046-01}{The French volume title is ``Histoire de Modèles''. ``Histoire'' can mean either a history or a story; ``History of Models'' preserves the singular title.} of Pursuing Stacks). At this point
the intuition emerges that $\infty$-groupoids must constitute
particularly adequate models for homotopy types,
n-groupoids corresponding to \uline{truncated} homotopy types
(with $\pi_i = 0$ for $i > n$). This same intuition, by very
different paths, was rediscovered by Ronnie Brown in Bangor and
some of his students, but using a rather restrictive notion of an $\infty$-groupoid
(which, among 1-connected homotopy types,
models only products of Eilenberg-Maclane spaces\frapparatus{V-FR-046-001}). It was
stimulated by a discursive correspondence with Ronnie Brown
that I finally resumed reflection, beginning with an attempt
to define a broader notion of an $\infty$-groupoid (subsequently renamed
``stack in $\infty$-groupoids'' or simply ``stack'', with the
understanding: over the punctual topos), which, one thing leading to another,
brought me to Pursuing Stacks. The volume ``History of Models''
moreover constitutes an entirely unforeseen digression
from the initial undertaking (the famous stacks being temporarily
forgotten, and not expected to reappear until around
page 1,000 ...).

\margin{46}
